% from ega4-18-fr.tex
\oldpage[IV-4]{221}

  de présentation finie. Pour que la $S$-immersion correspondante $j:Y\to X$ soit transversalement régulière relativement à $S$, il faut et il suffit qu'il existe $\lambda\geq\alpha$ tel que $j_\lambda:Y_\lambda\to X_\lambda$ soit transversalement régulière relativement à $S_\lambda$.
\end{enumerate}
\end{proposition}

\begin{enumerate}
  \item[(i)] La suffisance de la condition résulte de ($0$, 15.1.15). Pour prouver qu'elle est nécessaire, il suffira de montrer que tout point $x\in X$ possède un voisinage ouvert $V(x)$ pour lequel il existe un indice $\lambda=\lambda(x)\geq\alpha$ et un voisinage ouvert $V_\lambda$ de l'image $x_\lambda$ de $x$ dans $X_\lambda$ tels que $V(x)$ soit l'image réciproque de $V_\lambda$ et que la suite des restrictions à $V_\lambda$ des $f_{i\lambda}$ soit $(\sh F_\lambda|V_\lambda)$-transversalement régulière relativement à $S_\lambda$: en effet, il suffira alors d'appliquer (8.3.4) à l'ensemble ouvert $U_\lambda$ de $X_\lambda$, défini par la condition d'être le plus grand ouvert tel que les restrictions à $U_\lambda$ des $f_{i\lambda}$ forment une suite $(\sh F_\lambda|U_\lambda)$-transversalement régulière relativement à $S_\lambda$.

  Soient $s$ l'image de $x$ dans $S$, et pour tout $\lambda\geq\alpha$, soit $s_\lambda$ l'image de $s$ dans $S_\lambda$. Alors la fibre $X_s$ est limite projective des fibres $(X_\lambda)_{s_\lambda}$ et est localement noethérienne puisque le morphisme $X\to S$ est localement de présentation finie; comme en outre $(X_\alpha)_{s_\alpha}$ est quasi-compact et que les morphismes de transition $(X_\mu)_{s_\mu}\to(X_\lambda)_{s_\lambda}$ sont plats (puisqu'il en est ainsi de $\Spec(k(s_\mu))\to\Spec(k(s_\lambda))$) on peut appliquer aux sections $f_i\otimes1$ de $\sh F_s=\sh F\otimes_{\sh O_S}k(s)$ au-dessus de $X_s$ le résultat de (19.8.1, (i)); d'où l'existence d'un $\lambda\geq\alpha$ tel que les sections $f_{i\lambda}\otimes1$ de $(\sh F_\lambda)_{s_\lambda}=\sh F_\lambda\otimes_{\sh O_{S_\lambda}}k(s_\lambda)$ au-dessus de $(X_\lambda)_{s_\lambda}$ forment une suite $(\sh F_\lambda)_{s_\lambda}$-régulière. En outre (11.2.6), on peut supposer que $\sh F_\lambda$ est $S_\lambda$-plat au point $x_\lambda$. Il résulte alors de (11.3.8) qu'il existe un voisinage ouvert $V_\lambda$ de $x_\lambda$ dans $X_\lambda$ vérifiant les conditions requises.

  \item[(ii)] La suffisance de la condition résulte de (19.2.7, (ii)). Pour prouver sa nécessité, il suffit encore de montrer que tout point $x\in j(Y)$ admet un voisinage ouvert $V(x)$ pour lequel il existe un indice $\lambda=\lambda(x)\geq\alpha$ et un voisinage ouvert $V_\lambda$ de l'image $x_\lambda$ de $x$ dans $X_\lambda$ tels que $V(x)$ soit l'image réciproque de $V_\lambda$ et que la restriction $j_\lambda^{-1}(V_\lambda)\to V_\lambda$ de $j_\lambda$ soit transversalement régulière relativement à $S_\lambda$. Soient $s$ l'image de $x$ dans $S$, et pour tout $\lambda\geq\alpha$, soit $s_\lambda$ l'image de $s$ dans $S_\lambda$. Comme $X_s$ (resp. $Y_s$) est limite projective des fibres $(X_\lambda)_{s_\lambda}$ (resp. $(Y_\lambda)_{s_\lambda}$), on peut comme dans (i) utiliser (19.8.1, (ii)), et comme par hypothèse l'immersion $Y_s\to X_s$ est régulière, il existe un indice $\lambda(s)$ tel que pour $\lambda\geq\lambda(s)$ l'immersion $(Y_\lambda)_{s_\lambda}\to(X_\lambda)_{s_\lambda}$ soit régulière. En outre, en vertu de (19.2.4) il existe un voisinage ouvert quasi-compact $W(x)$ de $x$ dans $X$ tel que les morphismes structuraux $W(x)\to S$ et $Y\cap W(x)\to S$ soient plats; appliquant (11.2.6) et (8.2.11), on peut supposer qu'il existe un $\lambda\geq\lambda(s)$ et un voisinage ouvert quasi-compact $W(x_\lambda)$ de $x_\lambda$ dans $X_\lambda$, tels que $W(x)$ soit l'image réciproque de $W(x_\lambda)$ et que les restrictions à $W(x_\lambda)$ et à $Y_\lambda\cap W(x_\lambda)$ respectivement des morphismes structuraux $X_\lambda\to S_\lambda$ et $Y_\lambda\to S_\lambda$ soient plats. On conclut donc de (19.2.4) qu'il existe un voisinage ouvert $V_\lambda$ de $x_\lambda$ répondant à la question.
\end{enumerate}

\begin{remark}[19.8.3]
\label{IV.19.8.3}
Les démonstrations précédentes montrent que les énoncés de (19.8.1) et (19.8.2) subsistent lorsqu'on remplace les conditions concernant la régularité ou la régularité transversale dans tout $X$ (resp. $X_\lambda$) par ces conditions dans un voisinage d'un point $x\in X$ (resp. dans un voisinage de $x_\lambda$, projection de $x$).
\end{remark}

\subsection{\texorpdfstring{Suites $\sh F$-régulières et profondeur.}{Suites F-régulières et profondeur.}}
\EGAReaderTarget{ega4.c04.subsec.19-9}{subsection:IV.19.9}
\label{IV.19.9}

